% ============================================================
%  Capitolo 1 — I numeri
%  Matemagica — Libro di matematica per le scuole medie
%  Canton Ticino (DECS)
% ============================================================

\chapter{I numeri}

% ------------------------------------------------------------
\section*{Obiettivi di apprendimento}
% ------------------------------------------------------------

Al termine di questo capitolo l'allievo sarà in grado di:

\begin{itemize}
  \item Spiegare cosa sono i numeri e a cosa servono.
  \item Riconoscere e utilizzare le \textbf{cifre arabe} nel sistema di numerazione decimale.
  \item Leggere e scrivere semplici numeri con la \textbf{numerazione romana}.
  \item Identificare e collocare i \textbf{numeri naturali} sulla retta numerica.
  \item Comprendere il \textbf{sistema posizionale decimale}.
  \item Utilizzare correttamente i \textbf{simboli di confronto} tra numeri.
\end{itemize}

\bigskip

% ============================================================
\section{Cosa sono i numeri}
% ============================================================

L'idea di numero è una delle conquiste intellettuali più antiche dell'umanità.
Risale ad almeno 300\,000 anni fa, quando gli esseri umani vivevano nelle caverne
e dedicavano quasi tutte le proprie energie alla ricerca del cibo.
In quell'epoca non era ancora nota la capacità di contare:
l'idea di quantità era limitata ai piccoli insiemi che si potevano riconoscere
a colpo d'occhio — 1, 2, 3 e 4 oggetti.
Per quantità maggiori si ricorreva a espressioni generiche
come \emph{mucchio}, \emph{molti} o \emph{tanti}.

\medskip

Nel corso dei secoli, le civiltà hanno elaborato modalità diverse
per rappresentare i numeri, ciascuna con caratteristiche proprie.

\begin{definizione}
  Un \textbf{numero} è un insieme di simboli, chiamati \textbf{cifre},
  che indica una quantità.
\end{definizione}

% ------------------------------------------------------------
\subsection{Il sistema di numerazione arabo}
% ------------------------------------------------------------

Il sistema di numerazione che utilizziamo ogni giorno fu sviluppato
dalle popolazioni dell'\textbf{India} e successivamente diffuso in Europa
dagli \textbf{Arabi}.
Per questo motivo le nostre cifre sono chiamate \textbf{cifre arabe}
o \textbf{cifre indo-arabiche}.

\begin{regola}
  Il sistema arabo utilizza esattamente \textbf{dieci cifre}:
  \[
    0,\quad 1,\quad 2,\quad 3,\quad 4,\quad 5,\quad 6,\quad 7,\quad 8,\quad 9
  \]
  Con queste dieci cifre è possibile scrivere qualsiasi numero.
\end{regola}

\begin{esempio}
  Con le sole cifre $3$, $0$ e $7$ possiamo formare, tra gli altri,
  i numeri $307$, $370$, $703$ e $730$.
  Ogni combinazione rappresenta una quantità diversa.
\end{esempio}

% ------------------------------------------------------------
\subsection{Cenni di numerazione romana}
% ------------------------------------------------------------

Un altro sistema di numerazione importante nella storia è quello \textbf{romano},
usato nell'antica Roma e ancora oggi presente su orologi, monumenti e capitoli dei libri.

\begin{regola}
  La numerazione romana utilizza sette \textbf{simboli} fondamentali:

  \medskip
  \begin{center}
    \begin{tabular}{cc}
      \toprule
      \textbf{Simbolo} & \textbf{Valore} \\
      \midrule
      I   &   1 \\
      V   &   5 \\
      X   &  10 \\
      L   &  50 \\
      C   & 100 \\
      D   & 500 \\
      M   & 1000 \\
      \bottomrule
    \end{tabular}
  \end{center}

  \medskip
  I simboli si \textbf{sommano} se un simbolo di valore minore segue uno di valore maggiore
  (es.\ $\mathrm{VI} = 6$);
  si \textbf{sottraggono} se un simbolo di valore minore precede uno di valore maggiore
  (es.\ $\mathrm{IV} = 4$).
\end{regola}

\begin{esempio}
  \[
    \mathrm{VIII} = 8 \qquad
    \mathrm{XIV} = 14 \qquad
    \mathrm{XL} = 40 \qquad
    \mathrm{MMXXVI} = 2026
  \]
\end{esempio}

\begin{attenzione}
  \textbf{Errore comune — numerazione romana:}
  non è consentito ripetere lo stesso simbolo più di tre volte di fila.
  Quindi $40$ non si scrive $\mathrm{XXXX}$, ma $\mathrm{XL}$.
\end{attenzione}

\begin{sapevi}
  \textbf{Lo sapevi?}
  Il sistema romano non possiede un simbolo per lo zero.
  Questa mancanza rese i calcoli molto complessi.
  L'introduzione dello \textbf{zero} come cifra vera e propria,
  avvenuta in India attorno al V secolo d.C.,
  fu una delle rivoluzioni matematiche più importanti della storia.
\end{sapevi}

% ============================================================
\section{I numeri naturali}
% ============================================================

I numeri che fin dalle origini l'essere umano ha usato per \textbf{contare}
oggetti, animali, persone e giorni prendono il nome di \textbf{numeri naturali}.

\begin{definizione}
  I \textbf{numeri naturali} sono i numeri interi non negativi utilizzati per contare.
  Si indicano con il simbolo $\mathbb{N}$:
  \[
    \mathbb{N} = \{0,\; 1,\; 2,\; 3,\; 4,\; 5,\; \ldots\}
  \]
  I puntini ($\ldots$) indicano che la sequenza continua all'infinito.
\end{definizione}

% ------------------------------------------------------------
\subsection{La retta numerica}
% ------------------------------------------------------------

I numeri naturali possono essere rappresentati graficamente su una
\textbf{retta numerica}: una semiretta orientata su cui i numeri
sono collocati a distanze uguali tra loro.

\medskip

\begin{center}
\begin{tikzpicture}
  % Retta
  \draw[->, thick, color=blu] (-0.3,0) -- (8.5,0);
  % Tacche e numeri
  \foreach \x/\n in {0/0, 1/1, 2/2, 3/3, 4/4, 5/5, 6/6, 7/7} {
    \draw[thick] (\x,0.12) -- (\x,-0.12);
    \node[below=4pt] at (\x,0) {$\n$};
  }
  % Puntini
  \node at (8,0) {$\cdots$};
\end{tikzpicture}
\end{center}

\begin{regola}
  Sulla retta numerica:
  \begin{itemize}
    \item L'\textbf{origine} corrisponde al numero $0$.
    \item Procedendo verso destra, i numeri \textbf{aumentano}.
    \item Ogni numero naturale $n$ possiede un \textbf{successivo} $n + 1$
          e, ad eccezione dello zero, un \textbf{precedente} $n - 1$.
    \item Tra un numero e il suo successivo non esiste alcun altro numero naturale.
  \end{itemize}
\end{regola}

\begin{attenzione}
  \textbf{Attenzione — i numeri naturali sono infiniti:}
  non esiste un numero naturale più grande di tutti gli altri.
  Per indicare questo fatto si scrivono i puntini ($\ldots$)
  all'estremità destra della sequenza.
\end{attenzione}

% ============================================================
\section{Il sistema posizionale decimale}
% ============================================================

Il modo di rappresentare i numeri che utilizziamo abitualmente
si chiama \textbf{sistema posizionale decimale}.

\begin{definizione}
  Nel \textbf{sistema posizionale decimale}:
  \begin{itemize}
    \item Si utilizzano \textbf{dieci cifre}: $0, 1, 2, 3, 4, 5, 6, 7, 8, 9$.
    \item Il \textbf{valore} di ciascuna cifra dipende dalla
          \textbf{posizione} che occupa nel numero.
  \end{itemize}
\end{definizione}

\begin{esempio}
  Nel numero $\mathbf{352}$:
  \begin{itemize}
    \item la cifra $3$ si trova nella posizione delle \textbf{centinaia}
          e vale $3 \times 100 = 300$;
    \item la cifra $5$ si trova nella posizione delle \textbf{decine}
          e vale $5 \times 10 = 50$;
    \item la cifra $2$ si trova nella posizione delle \textbf{unità}
          e vale $2 \times 1 = 2$.
  \end{itemize}
  Quindi: $352 = 300 + 50 + 2$.
\end{esempio}

\begin{attenzione}
  \textbf{Errore comune — il ruolo dello zero:}
  la cifra $0$ non indica «nulla da scrivere» ma occupa una posizione precisa.
  Nel numero $305$, lo zero occupa la posizione delle decine
  e indica che non ci sono decine: $305 = 300 + 0 + 5$.
  Senza lo zero, $305$ e $35$ sarebbero indistinguibili.
\end{attenzione}

\begin{sapevi}
  \textbf{Lo sapevi?}
  Il termine \textbf{decimale} deriva dal latino \emph{decem}, che significa dieci.
  La scelta di dieci come base del sistema è probabilmente legata
  al fatto che l'essere umano possiede dieci dita,
  che in tempi antichi erano il primo strumento di calcolo.
\end{sapevi}

% ============================================================
\section{Simboli di confronto}
% ============================================================

In matematica si usano simboli specifici per confrontare due numeri
ed esprimere la loro relazione.

\begin{definizione}
  I \textbf{simboli di confronto} (o \textbf{simboli di relazione})
  indicano la relazione di ordine tra due numeri o espressioni.
\end{definizione}

\begin{regola}
  I simboli di confronto fondamentali sono:

  \medskip
  \begin{center}
    \begin{tabular}{cll}
      \toprule
      \textbf{Simbolo} & \textbf{Nome} & \textbf{Esempio} \\
      \midrule
      $>$   & maggiore di          & $7 > 3$ \quad (7 è maggiore di 3) \\
      $<$   & minore di            & $3 < 7$ \quad (3 è minore di 7) \\
      $\geq$ & maggiore o uguale a & $5 \geq 5$ \quad (5 è maggiore o uguale a 5) \\
      $\leq$ & minore o uguale a   & $4 \leq 6$ \quad (4 è minore o uguale a 6) \\
      $=$   & uguale a             & $2 + 3 = 5$ \\
      $\neq$ & diverso da          & $4 \neq 7$ \quad (4 è diverso da 7) \\
      \bottomrule
    \end{tabular}
  \end{center}
\end{regola}

\begin{esempio}
  \textbf{Uso dei simboli di confronto — passo per passo.}

  \medskip
  \textit{Confrontare $48$ e $52$.}

  \begin{enumerate}
    \item Osserviamo le cifre delle \textbf{decine}: $48$ ha la cifra $4$,
          mentre $52$ ha la cifra $5$.
    \item Poiché $4 < 5$, il numero $48$ è minore di $52$.
    \item Scriviamo: \quad $48 < 52$.
  \end{enumerate}
\end{esempio}

\begin{esempio}
  \textit{Verificare se $6 + 4 \neq 11$.}

  \begin{enumerate}
    \item Calcoliamo il primo membro: $6 + 4 = 10$.
    \item Confrontiamo con $11$: $10 \neq 11$ è \textbf{vero}.
  \end{enumerate}
\end{esempio}

\begin{attenzione}
  \textbf{Errore comune — direzione del simbolo:}
  il simbolo \emph{apre} sempre verso il numero maggiore.
  $5 > 3$ e $3 < 5$ esprimono la stessa relazione vista da due direzioni diverse.
  Un metodo utile: immaginare il simbolo come una freccia che punta
  verso il numero più piccolo.
\end{attenzione}

\begin{sapevi}
  \textbf{Lo sapevi?}
  I simboli $<$ e $>$ furono introdotti dal matematico inglese
  \textbf{Thomas Harriot} nel 1631, nell'opera \emph{Artis Analyticae Praxis}.
  Prima di allora, i matematici esprimevano le relazioni di ordine
  soltanto a parole.
\end{sapevi}

% ============================================================
\section*{Riepilogo del capitolo}
% ============================================================

\begin{itemize}
  \item Un \textbf{numero} è un insieme di simboli (cifre) che indica una quantità.
  \item Il \textbf{sistema arabo} usa dieci cifre ($0$–$9$);
        il \textbf{sistema romano} usa sette simboli ($\mathrm{I, V, X, L, C, D, M}$).
  \item I \textbf{numeri naturali} ($\mathbb{N} = \{0, 1, 2, 3, \ldots\}$)
        servono a contare; sono infiniti e si rappresentano sulla retta numerica.
  \item Nel \textbf{sistema posizionale decimale} il valore di una cifra
        dipende dalla posizione che occupa nel numero.
  \item I \textbf{simboli di confronto}
        ($>, <, \geq, \leq, =, \neq$)
        esprimono la relazione tra due quantità.
\end{itemize}
